\chapter{限界と今後の課題}
\section{限界}
\begin{enumerate}[leftmargin=2em]
\item 第3回学習は総20 epoch計画のepoch 5で停止し、epoch 6--20は未実施である。SiLog基準runの収束・過学習は未確定である。
\item 無料Colab GPUの利用制限により追加学習は停止した。追加課金や制限回避は行っていない。
\item validationはDay047の128画像に限られ、未知撮影日・未知圃場への汎化は未検証である。
\item 複数seedによるばらつきは未計測であり、画像単位信頼区間や有意差も報告していない。
\item independent test319組は未取得・未使用、過去52組も学習・モデル選択には未使用である。
\item GTは写真測量由来カメラZ-depthで地上実測値と同一ではない。地上垂直草丈GTとの比較はない。
\item 稲UAV画像および稲の垂直草丈GTは未使用であり、稲高精度を実証していない。
\item 入力518 pxへの縮小は微細構造を弱める可能性があるが、高解像度tileとの比較はしていない。
\end{enumerate}
\section{次の検証候補}
第一に、保存済みD checkpointから同じ20 epoch scheduler、同一署名コード・分割を保ってepoch 6--20を再開する。epoch毎に絶対深度と局所形状を記録する。epoch 5でRMSEが改善中であることから残りepochを完了すれば曲線の把握に価値はあるが、改善の継続は保証されない。
第二に、SiLog基準runの評価後、新規別runで局所差分損失を比較する。全画面MAE/RMSEを保持しながら、GT階級別振幅・符号・複数間隔を事前に定義する。基準runへ損失変更を混ぜない。第三に重複tile、高解像度入力、マルチスケール推論を比較する。GPU前に無料GPU割当とVRAMを再計測する。第四にAの端要因をpadding、crop位置・アスペクト比、端損失寄与など個別変数として切り分けるが、端除外を主指標にしない。第五に稲画像、地面・頂部、カメラ姿勢と地上実測高さを収集し、Z-depthから草丈へ変換する方法を検証する。モデル選択完了後にindependent testを最終評価する。